\section{Related Work}
\label{sec:related}

Our framework draws on and departs from three research traditions: cultural
evolution theory, memetics, and computational models of opinion dynamics.

\subsection{Cultural Evolution}

The mathematical study of cultural transmission begins with
\citet{cavalli1981cultural}, who adapted population genetics models to
cultural traits, distinguishing vertical (parent--child), horizontal
(peer--peer), and oblique (elder--youth) transmission modes.
\citet{boyd2005origin} extended this program with models of biased
transmission, conformist bias, and prestige bias, showing that cultural
evolution can produce group-level adaptations not reducible to individual
selection. \citet{mesoudi2011cultural} provides a comprehensive synthesis.

These models treat cultural traits as discrete variants propagating through
a population --- formally, alleles in a cultural genome. The framework is
powerful for small-scale traits (tool designs, food taboos, phonemes) but
encounters difficulties with \emph{ideologies}: large, internally structured
systems of belief that interact with institutions, respond to economic
conditions, and actively reshape the environments in which they propagate.
The cultural evolution tradition typically assumes a fixed or slowly varying
selective environment; our framework makes the environment endogenous.

\citet{bisin2001economics} developed an economic model of cultural
transmission where parents invest in socializing children into their own
cultural group. Their model produces long-run coexistence of cultural types
through a mechanism analogous to our latent sympathy: minority groups invest
more heavily in transmission, preventing assimilation. Our latent variable
$y_i$ generalizes this insight beyond the parent--child dyad.

\subsection{Memetics}

\citet{dawkins1976selfish} proposed that ideas (``memes'') replicate, mutate,
and compete in a manner analogous to genes. The analogy has been productive
as a heuristic but has resisted formalization. \citet{shennan2002genes}
applied memetic reasoning to archaeological data with some success, but the
field has struggled with a core problem: unlike genes, ideas do not have a
well-defined unit of replication. A political ideology is not a single meme
but a complex of mutually reinforcing propositions, aesthetic preferences,
institutional practices, and identity markers.

Our framework sidesteps this problem by working at the level of entire
ideological systems rather than individual memes. We model ideologies as
indivisible competitors on a simplex --- each characterized by a vector of
responses to environmental conditions --- rather than attempting to decompose
them into memetic atoms. This is a deliberate coarse-graining: it sacrifices
internal structure for tractability at the population level.

\subsection{Opinion Dynamics and Social Contagion}

A large literature in computational social science models opinion spread as a
contagion or threshold process. \citet{granovetter1978threshold} introduced
threshold models where individuals adopt a behavior when a sufficient
fraction of their neighbors have adopted it. \citet{watts1998collective}
showed that small-world network structure can produce cascading adoption.
\citet{acemoglu2011bayesian} studied Bayesian learning in social networks,
showing how network structure determines whether populations converge to
correct beliefs.

These models typically operate at the individual level and model opinion as a
continuous variable or binary choice. Our framework operates at the
population level and models competition between multiple discrete ideologies.
The two approaches are complementary: network models explain the
micro-mechanisms of spread (why \emph{this} person adopted \emph{this} idea
now), while our framework explains the macro-dynamics (why liberalism grew
in the 19th century and contracted in the 21st).

\subsection{Empirical Studies of Ideological Change}

The empirical literature provides the calibration targets for our model.
\citet{funke2016going} analyzed 800 elections in 20 democracies over 140
years and found that financial crises increase far-right vote share by
approximately 30\%. \citet{inglehart1997modernization} documented the
``post-materialist'' value shift across generations using 50 years of World
Values Survey data, confirming the generational transmission mechanism.
\citet{jennings1968transmission} showed that parental party identification
predicts children's at $r = 0.47$--$0.59$. \citet{pierson2000increasing}
formalized the ``policy ratchet'' --- the tendency of institutional changes
to become self-reinforcing --- providing the empirical basis for our
institutional amplification mechanism. \citet{norris2019cultural}
documented the rise of authoritarian populism as a cultural backlash against
post-materialist values, a dynamic our model captures through the
interaction between environmental variables and ideological fitness.

\citet{king2008ordinary} showed that unemployment correlated with radical
vote share across Weimar Germany's electoral districts, providing perhaps the
cleanest quantitative evidence for crisis recruitment. These findings ground
our treatment of exogenous shocks as temporary increases in the
susceptibility of populations to ideological change.

\subsection{What Is Missing}

No existing framework, to our knowledge, combines all three features we
address: dormancy with martyrdom dynamics, endogenous environment
modification, and hybridization as speciation. Cultural evolution models
handle transmission but not institutional feedback. Memetic models handle
competition but not latency. Opinion dynamics models handle micro-level
spread but not macro-level ideological competition. Our contribution is to
integrate these mechanisms into a single, mathematically tractable system.
